\chapter{Experimental Setup}
\label{ch:setup}

\section{Data}
\begin{itemize}
  \item \textbf{2D cells (scattering).} The first samples of the InverseBench \texttt{inv-scatter-test} set:
  $128\times128$ permittivity images of 1--6 HL60 nuclei simulated with CytoPacq~\citep{wiesner2019cytopacq}, values in
  $[0,1]$.
  \item \textbf{3D cells (scattering).} (a) Synthetic HL60-like nuclei generated by our code: 1--6 soft, rotated
  ellipsoids with smooth internal texture in a $32\times128\times128$ volume. (b) Real CytoPacq nuclei from BBBC024
  of the Broad Bioimage Benchmark Collection~\citep{ljosa2012bbbc} (ground-truth label masks, 20 nuclei per volume),
  converted to foreground masks, cropped to the central half in $x$ and $y$ and resampled to
  $32\times128\times128$.
  \item \textbf{Fluid (4D).} Initial vorticity fields from the InverseBench \texttt{navier-stokes-test} set
  ($Re=200$), rolled forward with the InverseBench pseudo-spectral solver to 21 frames spaced $\Delta t=0.05$
  apart, divided by 10 to lie in about $[-1,1]$. Frames 2, 6, 10, 14 and 18 are held out; the rest are observed
  completely or at a random 10\% of pixels per frame.
\end{itemize}

\section{Forward models and noise}
Scattering uses the InverseBench configuration: domain $0.18\times0.18$\,m, wavelength six grid cells, 20
transmitters and 360, 180 or 60 receivers on a circle of radius 1.6\,m, complex Gaussian noise $\sigma=10^{-4}$ with a
fixed seed per sample so that Gaussian and baseline fits see identical measurements.

\section{Hyper-parameters}
\begin{table}[h]
  \centering\small
  \caption{Default settings (overridable from the command line; every run stores its exact settings).}
  \label{tab:hparams}
  \resizebox{\linewidth}{!}{%
\begin{tabular}{lll}
    \toprule
    Setting & Scattering 2D / 3D & Fluid 4D \\
    \midrule
    Iterations & 3000 & 3000 \\
    Optimiser, learning rate & Adam, 0.05 (Gaussians), 0.01 (grid) & Adam, 0.01 (Gaussians), $10^{-3}$ (network) \\
    Schedule & cosine & cosine \\
    Max.\ / initial Gaussians & 2000 / 250 (2D); 8000 / 2000 (3D) & 2000 / 2000 \\
    Window half-size $h$ & 8 (2D), 6 (3D) & 10, periodic \\
    Densification & every 250 it., until 60\%, $+n_{\mathrm{init}}/4$ & none \\
    Sparsity $\lambda$ & 0, 0.1, 0.3, 1, 3 & -- \\
    Temporal smoothness $\lambda_T$ & -- & 0 (initial), 1 (final) \\
    Deformation & -- & MLP (3$\times$128), Mamba and bidirectional Mamba (2 blocks, $d=64$, state 16) \\
    Batch & one sample & 4 frames \\
    Samples per setting & 3 (5 in the final 2D runs) & 3 sequences \\
    \bottomrule
  \end{tabular}}
\end{table}

\section{Hardware and software}
All experiments ran on a single NVIDIA Tesla T4 (16\,GB) on Google Colab or Kaggle, with Python 3, PyTorch, NumPy,
SciPy and scikit-image. A 2D scattering fit takes about 13\,s; a 3D fit a few minutes; a 4D fit with the MLP about
3\,min and with Mamba 15--28\,min (sequential scan; the longer time with temporal smoothness).

\section{Protocol}
For each sample, receiver count and method, measurements are generated once and both the Gaussian and the grid
method are fitted to them. Reported numbers are means over samples; per-sample values, figures and run
configurations are stored in the repository (\texttt{results/}).
